% === Begin Label Data ===
% ID: 255
% 难度星级: 
% 标签: 
% 备注: 
% 组卷引用次数: 0
% === End  Label Data ===

\begin{problem}{2026}{M}{江苏南京市、盐城市2026届高三年级第一次模拟考试}{7}{三角函数}
设 $\sigma$ 和 $\tau$ 表示坐标平面内的几何变换，$\sigma$ 表示将几何对象绕原点 $O$ 逆时针旋转 $\dfrac{\pi}{12}$，$\tau$ 表示几何对象关于 $y$ 轴对称，$\sigma^k$（$k\in\mathbb{N}^*$）表示连续 $k$ 次 $\sigma$ 变换．已知角 $\alpha$ 的终边经过点 $(-2,1)$，若对角 $\alpha$ 的终边先进行 $\tau$ 变换，再进行 $\sigma^3$ 变换，得到角 $\beta$ 的终边，则 $\tan\beta=$
(\hspace{1cm})
\begin{choices}
\choice{{$-3$}}
\choice{{ $-\dfrac{1}{3}$}}
\choice{{ $\dfrac{1}{3}$}}
\choice{{$3$}}
\end{choices}
\end{problem}